\chapter{Reinforcement Learning Basics}
\label{ch:rl}

\chapterguide%
  {Probability, expected value, optimization, and the idea of sequential decisions.}
  {define states, actions, rewards, policies, and values; apply Q-learning and evaluate separately from exploration.}

Reinforcement learning (RL) learns sequential decisions from rewards rather than correct action labels. Actions change future states, so the objective is long-term return.

\begin{coursechoice}
Lab~6 implements tabular RL. Neural function approximation for deep RL lies beyond this course.
\end{coursechoice}

\section{The agent-environment loop}

At time $t$:

\begin{mlsteps}
  \item The agent observes a state $s_t$.
  \item It chooses an action $a_t$.
  \item The environment returns a reward $r_{t+1}$ and a new state $s_{t+1}$.
  \item The agent updates what it believes about the action.
\end{mlsteps}

\begin{mlfigure}
\begin{tikzpicture}[
  node distance=3.5cm,
  box/.style={draw=MLBlue,fill=MLPaleBlue,rounded corners=2mm,minimum width=3cm,minimum height=1.2cm,align=center,font=\bfseries\color{MLNavy}},
  arr/.style={-{Stealth[length=2.4mm]},thick,draw=MLTeal}
]
\node[box] (agent) {Agent\\chooses an action};
\node[box,right=of agent] (env) {Environment\\changes state};
\draw[arr] (agent.north east) to[bend left=18] node[above,font=\small]{action $a_t$} (env.north west);
\draw[arr] (env.south west) to[bend left=18] node[below,font=\small]{state $s_{t+1}$, reward $r_{t+1}$} (agent.south east);
\end{tikzpicture}
\end{mlfigure}

\begin{workedexample}
A warehouse robot tracks location and battery, chooses movement or charging, and balances time and collision costs against delivery rewards.
\end{workedexample}

\section{A simpler case: multi-armed bandits}

A bandit chooses among uncertain-reward actions without changing state, balancing learning with reward.

This creates the central RL problem:

\begin{description}[leftmargin=1.35in,style=nextline]
  \item[Exploration] Try uncertain actions to learn more about them.
  \item[Exploitation] Choose the action currently believed to be best.
\end{description}

An $\varepsilon$-greedy strategy explores with probability $\varepsilon$ and otherwise exploits.

\section{Markov decision processes}

A classical RL problem is often written as a Markov decision process (MDP) with:

\begin{itemize}
  \item a set of states $\mathcal{S}$;
  \item a set of actions $\mathcal{A}$;
  \item transition probabilities $P(s'\given s,a)$;
  \item a reward rule;
  \item a discount factor $\gamma\in[0,1)$.
\end{itemize}

Under the Markov assumption, the current state and action contain the information needed for the next-state distribution.

\section{Return and discounting}

The return from time $t$ is the discounted sum of future rewards:
\[
  G_t=r_{t+1}+\gamma r_{t+2}+\gamma^2r_{t+3}+\cdots.
\]

The discount factor $\gamma$ controls how much future rewards matter.

\begin{itemize}
  \item Small $\gamma$: favor immediate rewards.
  \item Large $\gamma$: value long-term rewards more strongly.
\end{itemize}

\begin{workedexample}
Suppose future rewards are $5$, $4$, and $3$, with $\gamma=0.9$. Then
\[
G_t=5+0.9(4)+0.9^2(3)=11.03.
\]
The later rewards count, but each is discounted.
\end{workedexample}

\section{Policies and value functions}

A \term{policy} $\pi(a\given s)$ describes how the agent chooses actions in each state.

The state-value function is
\[
  V^{\pi}(s)=\mathbb{E}_{\pi}[G_t\given s_t=s],
\]
the expected return from state $s$ while following policy $\pi$.

The action-value function is
\[
  Q^{\pi}(s,a)=\mathbb{E}_{\pi}[G_t\given s_t=s,a_t=a],
\]
the expected return from taking action $a$ in state $s$ and then following $\pi$.

\begin{intuition}
$V(s)$ asks, ``How good is this situation?'' $Q(s,a)$ asks, ``How good is this particular move from this situation?''
\end{intuition}

\section{The Bellman idea}

A long-term value can be broken into two parts:

\[
\text{value now}=\text{reward next}+\text{discounted value later}.
\]

For an optimal action-value function,
\[
Q^*(s,a)=\mathbb{E}\left[r_{t+1}+\gamma\max_{a'}Q^*(s_{t+1},a')\mid s_t=s,a_t=a\right].
\]
This recursive relationship is the basis of Q-learning.

\section{Tabular Q-learning}

When the state and action spaces are small, store one Q-value for every state-action pair in a table.

After observing transition $(s_t,a_t,r_{t+1},s_{t+1})$, update
\[
Q(s_t,a_t)\leftarrow Q(s_t,a_t)+\alpha\left[r_{t+1}+\gamma\max_{a'}Q(s_{t+1},a')-Q(s_t,a_t)\right].
\]

The bracketed quantity is the \term{temporal-difference error}: the difference between the new backup value and the current estimate.

\begin{mlsteps}
  \item Initialize the Q-table, often with zeros.
  \item Observe the current state.
  \item Choose an action using an exploration strategy such as $\varepsilon$-greedy.
  \item Perform the action and observe reward and next state.
  \item Update the chosen Q-table entry.
  \item Repeat over many episodes, gradually reducing exploration when appropriate.
\end{mlsteps}

Lab~6 trains on a $4\times4$ grid. After 5,000 episodes, its greedy policy reaches the goal in six moves for return $0.8$; the saved model is the Q-table.

\begin{workedexample}
With $Q(s,a)=2$, $r=1$, $\gamma=0.9$, next-state maximum 4, and $\alpha=0.5$, the backup is $4.6$ and TD error $2.6$:
\[
Q(s,a)\leftarrow2+0.5(2.6)=3.3.
\]
The estimate moves halfway toward the new backup value.
\end{workedexample}

\section{One Q-learning update, drawn}

The diagram separates a Q-update into its current estimate, backup value, and step size.

\begin{mlfigure}
\begin{tikzpicture}[
  node distance=0.42cm and 0.8cm,
  bx/.style={draw=MLBlue,fill=MLPaleBlue,rounded corners=1.5mm,minimum width=2.5cm,minimum height=0.85cm,align=center,font=\scriptsize\bfseries\color{MLNavy}},
  gd/.style={draw=MLGreen,fill=MLPaleTeal,rounded corners=1.5mm,minimum width=2.5cm,minimum height=0.85cm,align=center,font=\scriptsize\bfseries\color{MLNavy}},
  op/.style={draw=MLOrange,fill=MLPaleOrange,rounded corners=1.5mm,minimum width=3.6cm,minimum height=0.82cm,align=center,font=\scriptsize\bfseries\color{MLNavy}},
  arr/.style={-{Stealth[length=2mm]},thick,draw=MLTeal}
]
\node[bx] (old) at (-2.35,1.55) {Old estimate\\$Q(s,a)$};
\node[gd] (backup) at (2.35,1.55) {Backup value\\$r+\gamma\max_{a'}Q(s',a')$};
\node[op] (update) at (0,0.30)
  {Move a fraction $\alpha$ toward the backup value\\$Q+\alpha(\text{backup}-Q)$};
\node[bx] (new) at (0,-1.05) {New estimate\\$Q(s,a)$};
\draw[arr] (old.south)--(update.north west);
\draw[arr] (backup.south)--(update.north east);
\draw[arr] (update.south)--(new.north);
\node[font=\scriptsize\itshape\color{MLMuted},below=0.25cm of new,align=center]
  {what actually happened\\nudges what you believed};
\end{tikzpicture}
\end{mlfigure}

\section{Exploration strategies}

\begin{mltable}
\begin{tabularx}{\linewidth}{@{}>{\bfseries\leavevmode\color{MLNavy}}p{0.27\linewidth}X@{}}
\toprule
\rowcolor{MLPaleBlue!92}
Exploration strategy & How it encourages exploration \\
\midrule
$\varepsilon$-greedy & Usually chooses the best-known action, but chooses a random action with probability $\varepsilon$. \\
Optimistic initialization & Starts Q-values high so untried actions look attractive. \\
Upper confidence methods & Prefer actions with either strong estimated reward or high uncertainty. \\
\bottomrule
\end{tabularx}
\end{mltable}

Too little exploration can trap the agent in a poor policy. Too much exploration wastes reward and may be unsafe.

\section{Evaluation and responsible use}

Policies change the data they encounter. Evaluate across seeds and report mean return and variation.

Important concerns include:

\begin{itemize}
  \item reward design that accidentally encourages unwanted behavior;
  \item unsafe exploration in physical, medical, financial, or public systems;
  \item changing environments that make past experience unreliable;
  \item delayed rewards that make credit assignment difficult;
  \item simulation assumptions that do not match the real world.
\end{itemize}

For high-stakes applications, exploration should occur in a safe simulator or under strict constraints before real-world use.

\begin{modelcard}{Classical Reinforcement Learning - Quick Guide}
\begin{tabularx}{\linewidth}{@{}>{\bfseries\leavevmode\color{MLNavy}}p{0.14\linewidth}X@{}}
Best when & Actions affect future states and rewards, and interaction or a realistic simulator is available. \cr
Start with & A bandit or small tabular MDP before moving to larger problems. \cr
Main controls & Learning rate $\alpha$, discount factor $\gamma$, exploration rate $\varepsilon$, state representation, reward design. \cr
Main risks & Unsafe exploration, reward hacking, unstable estimates, poor state definition, unrealistic simulation. \cr
\end{tabularx}
\end{modelcard}

\begin{keyidea}
Reinforcement learning does not learn from a correct answer for every action. It learns from consequences, balancing immediate reward, future reward, and the need to explore.
\end{keyidea}

\section{Bandit exploration behavior}

Three bandit arms pay $0.2$, $0.5$, and $0.8$ on average. An $\varepsilon=0.1$ agent over 1,000 rounds might produce:

\begin{mltable}
\begin{tabularx}{0.75\linewidth}{@{}>{\bfseries\leavevmode\color{MLNavy}}p{0.22\linewidth}X X@{}}
\toprule
\rowcolor{MLPaleBlue!92}
Machine & True mean reward & Pulls in 1,000 rounds \\
\midrule
A & $0.2$ & about 40 \\
B & $0.5$ & about 40 \\
C & $0.8$ & about 920 \\
\bottomrule
\end{tabularx}
\end{mltable}

\begin{keyidea}
The agent mostly exploits the best arm while exploring in one decision out of ten. Exploration reduces premature commitment at the cost of immediate reward.
\end{keyidea}

\begin{mlfigure}
\begin{tikzpicture}
\begin{axis}[
  width=0.66\linewidth, height=4.3cm, axis lines=left,
  title style={font=\small\bfseries\color{MLNavy}}, title={Decaying $\varepsilon$: curious early, decisive later},
  xlabel={Episode}, ylabel={$\varepsilon$},
  label style={font=\scriptsize\color{MLMuted}}, tick label style={font=\scriptsize\color{MLMuted}},
  domain=0:1000, samples=100, xmin=0, xmax=1000, ymin=0, ymax=1.05,
  legend style={font=\scriptsize,draw=none,fill=none,at={(0.98,0.96)},anchor=north east}
]
\addplot[MLBlue,very thick]{max(0.05, exp(-x/250))};   \addlegendentry{exponential decay}
\addplot[MLTeal,very thick,dashed]{max(0.05, 1 - x/800)}; \addlegendentry{linear decay}
\addplot[MLMuted,thick,dotted]{0.1};                   \addlegendentry{fixed $\varepsilon=0.1$}
\end{axis}
\end{tikzpicture}
\end{mlfigure}

\section{Tabular Q-learning dynamics}

\begin{mltable}
\begin{tabularx}{\linewidth}{@{}>{\bfseries\leavevmode\color{MLNavy}\ttfamily}p{0.15\linewidth}p{0.16\linewidth}X@{}}
\toprule
\rowcolor{MLPaleBlue!92}
\normalfont\bfseries\leavevmode\color{MLNavy}Symbol & \normalfont\bfseries\leavevmode\color{MLNavy}Typical value & \normalfont\bfseries\leavevmode\color{MLNavy}What raising it does \\
\midrule
alpha & 0.01--0.2 & Learn faster from each experience, but more erratically \\
gamma & 0.9--0.99 & Care more about distant future rewards \\
epsilon & 1.0 down to 0.05 & Explore more; too high and the agent never commits \\
episodes & thousands & More practice; tabular methods need a lot \\
\bottomrule
\end{tabularx}
\end{mltable}

\begin{pitfall}
Continuing tasks generally use $\gamma<1$ for a finite discounted return; finite episodes can allow $\gamma=1$. Fixed exploration can help in changing environments but costs reward. Evaluate exploratory and greedy policies separately.
\end{pitfall}

\begin{libnote}
\lib{gymnasium} (the maintained successor to OpenAI Gym) provides ready-made environments with a uniform interface, so you can practice without writing the environment mechanics yourself. For problems beyond small tables --- for example, image observations or continuous state spaces --- \lib{stable-baselines3} provides implementations of standard deep-RL algorithms on top of PyTorch.
\end{libnote}

\begin{keyidea}
Large or continuous state spaces require function approximation; neural approximators lead to deep RL.
\end{keyidea}

\labsection{6}{Reinforcement learning}{sec:lab06-reinforcement}

\begin{labproject}{Learn a policy by interacting, then save it}
\textbf{Goal.} Train tabular Q-learning, evaluate the greedy policy, visualize values/actions, and save the Q-table.

\medskip\noindent\textbf{Setup.}\par
\begin{lstlisting}[style=mlpython]
from pathlib import Path

import matplotlib.pyplot as plt
import numpy as np

PROJECT_ROOT = Path.cwd().parent
MODELS = PROJECT_ROOT / "models"
MODELS.mkdir(exist_ok=True)
\end{lstlisting}

\medskip\noindent\textbf{Part A --- Tabular Q-learning in a 4$\times$4 GridWorld.}\par

Use a $4\times4$ grid: start top-left, goal bottom-right, and a trap. Ordinary moves cost $-0.04$, the goal pays $+1$, and the trap pays $-1$ and terminates. Walls leave position unchanged. \texttt{step} returns state, reward, and termination.

\begin{lstlisting}[style=mlpython]
rng = np.random.default_rng(42)

ROWS, COLS = 4, 4
START, GOAL, TRAP = (0, 0), (3, 3), (1, 3)
ACTIONS = [(-1, 0), (1, 0), (0, -1), (0, 1)]        # up, down, left, right
ACTION_NAMES = ["up", "down", "left", "right"]
ARROWS = {"up": (0, 1), "down": (0, -1), "left": (-1, 0), "right": (1, 0)}

def step(state, action):
    """Apply one action. Returns (next_state, reward, done)."""
    r, c = state
    dr, dc = ACTIONS[action]
    next_state = (min(max(r + dr, 0), ROWS - 1), min(max(c + dc, 0), COLS - 1))
    if next_state == GOAL:
        return next_state, 1.0, True
    if next_state == TRAP:
        return next_state, -1.0, True
    return next_state, -0.04, False
\end{lstlisting}

Initialize state--action Q-values to zero. Update toward reward plus discounted best next-state value, with step size $\alpha$ and discount $\gamma$. Decay exploration from $\varepsilon=1$.

\begin{lstlisting}[style=mlpython]
Q = np.zeros((ROWS, COLS, len(ACTIONS)))
alpha, gamma = 0.15, 0.95
epsilon, min_epsilon = 1.0, 0.05

def choose_action(state, explore=True):
    if explore and rng.random() < epsilon:
        return int(rng.integers(len(ACTIONS)))       # explore: random move
    return int(np.argmax(Q[state]))                  # exploit: best known move

training_returns = []
for episode in range(5000):
    state, episode_return = START, 0.0
    for _ in range(100):
        action = choose_action(state)
        next_state, reward, done = step(state, action)
        episode_return += reward
        target = reward if done else reward + gamma * np.max(Q[next_state])
        Q[state][action] += alpha * (target - Q[state][action])
        state = next_state
        if done:
            break
    training_returns.append(episode_return)
    epsilon = max(min_epsilon, epsilon * 0.9985)
\end{lstlisting}

Plot returns smoothed over 100 episodes as exploration decays.

\begin{lstlisting}[style=mlpython]
window = 100
moving = np.convolve(training_returns, np.ones(window) / window, mode="valid")
plt.figure(figsize=(6, 3.6))
plt.plot(np.arange(window - 1, len(training_returns)), moving)
plt.axhline(0.8, linestyle="--", color="grey", label="best possible return (6 moves)")
plt.xlabel("Episode")
plt.ylabel("Mean return over last 100 episodes")
plt.title("Q-learning training curve")
plt.legend()
plt.tight_layout()
plt.show()
\end{lstlisting}

Evaluate separately with exploration off, choosing the largest Q-value in each cell.

\begin{lstlisting}[style=mlpython]
successes, returns = 0, []
for _ in range(200):
    state, total = START, 0.0
    for _ in range(100):
        state, reward, done = step(state, choose_action(state, explore=False))
        total += reward
        if done:
            successes += int(state == GOAL)
            break
    returns.append(total)

print("success rate:", successes / 200)
print("mean return:", round(float(np.mean(returns)), 3))
print("learned action at start:", ACTION_NAMES[choose_action(START, explore=False)])
\end{lstlisting}

\textbf{Inspect.} Six moves yield $5(-0.04)+1=0.80$. A 100\% evaluated success rate at return 0.8 indicates the shortest route.

\medskip\noindent\textbf{Part B --- The Q-table is the model: save it and read the policy.}\par

Save the Q-table; deployment chooses each state’s highest-valued action without retraining.

\begin{lstlisting}[style=mlpython]
q_path = MODELS / "gridworld_q_table.npy"
np.save(q_path, Q)

Q_loaded = np.load(q_path)
policy = np.full((ROWS, COLS), "", dtype=object)
for r in range(ROWS):
    for c in range(COLS):
        policy[r, c] = "GOAL" if (r, c) == GOAL else "TRAP" if (r, c) == TRAP \
            else ACTION_NAMES[int(np.argmax(Q_loaded[r, c]))]
print(policy)
\end{lstlisting}

Plot greedy actions as arrows over maximum Q-values.

\begin{lstlisting}[style=mlpython]
def plot_policy(Q_table, title):
    values = Q_table.max(axis=2)
    plt.figure(figsize=(4.4, 4.2))
    plt.imshow(values, cmap="YlGn")
    plt.colorbar(label="value of the best action")
    for r in range(ROWS):
        for c in range(COLS):
            if (r, c) == GOAL:
                plt.text(c, r, "GOAL", ha="center", va="center", fontweight="bold")
            elif (r, c) == TRAP:
                plt.text(c, r, "TRAP", ha="center", va="center", color="red", fontweight="bold")
            else:
                dx, dy = ARROWS[ACTION_NAMES[int(np.argmax(Q_table[r, c]))]]
                plt.arrow(c, r, 0.3 * dx, -0.3 * dy, head_width=0.15, color="black")
    plt.xticks(range(COLS))
    plt.yticks(range(ROWS))
    plt.title(title)
    plt.tight_layout()
    plt.show()

plot_policy(Q_loaded, "Greedy policy and state values")
\end{lstlisting}

\textbf{Inspect.} Trace routes toward the goal and around the trap; compare neighboring state values.

\begin{tryit}
Change the step cost, trap, exploration schedule, or discount. Predict the policy change, retrain, and check with \texttt{plot\_policy}.
\end{tryit}
\end{labproject}


\begin{quickcheck}
\begin{enumerate}
  \item What are the state, action, reward, and transition in a reinforcement-learning formulation?
  \item What two pieces of information appear in the tabular Q-learning backup value?
  \item Why should the exploratory training policy be evaluated separately from the final greedy policy?
\end{enumerate}
\end{quickcheck}
